% Generated draft: requires mathematical and source review.
\documentclass[11pt]{article}
\usepackage[margin=25mm]{geometry}
\usepackage{fontspec}
\setmainfont{Latin Modern Roman}
\usepackage{amsmath,amssymb,mathtools,unicode-math}
\setmathfont{Latin Modern Math}
\usepackage{enumitem,needspace,xurl,hyperref}
\setlength{\parindent}{0pt}
\setlength{\parskip}{5pt}
\setlength{\emergencystretch}{3em}
\newcommand{\N}{\mathbb{N}}
\newcommand{\Z}{\mathbb{Z}}
\newcommand{\Q}{\mathbb{Q}}
\newcommand{\R}{\mathbb{R}}
\newcommand{\C}{\mathbb{C}}
\newcommand{\F}{\mathbb{F}}
\newcommand{\A}{\mathbb{A}}
\newcommand{\PP}{\mathbb{P}}
\newcommand{\E}{\mathbb{E}}
\newcommand{\eps}{\varepsilon}
\newcommand{\dd}{\,\mathrm{d}}
\begin{document}
\section*{UnsolvedMath ID: 30000488}
\textbf{Source problem code:} catalogue entry
\subsection{Definitions and mathematical statement}
\textbf{Structure problem.} Let \(p\) be a prime, \(F\) a number field, and \(L/F\) a pro-\(p\) Galois extension unramified outside a finite set \(S\) containing the primes above \(p\). Determine the Galois group \(G_L\) of the maximal unramified pro-\(p\) extension of \(L\), and its maximal abelian quotient \(X_L\). A pro-\(p\) group is a profinite inverse limit of finite \(p\)-groups, so the hypothesis permits finite as well as infinite \(L/F\); the title's reference to infinite extensions does not add an explicit infinitude assumption.\par Fixing a compatible separable closure and the intended notion of unramifiedness, write \(L^{\mathrm{ur},p}\) for the compositum of the unramified finite Galois \(p\)-extensions of \(L\). Then\[G_L=\operatorname{Gal}(L^{\mathrm{ur},p}/L),\qquad X_L=G_L/\overline{[G_L,G_L]}.\] The closure of the commutator subgroup is essential: this is the maximal continuous abelian pro-\(p\) quotient. Determining \(X_L\) alone does not determine the full nonabelian \(G_L\). The condition on ramification outside \(S\) applies to the given extension \(L/F\), while \(L^{\mathrm{ur},p}/L\) is requested to be unramified in the chosen sense; these conditions are not interchanged.\par The source does not specify whether archimedean ramification is allowed, a distinction particularly relevant when \(p=2\), or the valuation-theoretic convention over an infinite algebraic base \(L\). It also does not fix a family of towers, provide arithmetic data determining one example, or specify the desired form of a general answer, such as presentations, finiteness criteria, or invariants. Thus this is a structural research program over the stated class, not one numerical group computation. No \(\Z_p\)-extension, cyclotomic, or finite-generation hypothesis is added.\par The supplied triage checked 2026-08-21 labels the problem partially solved. It says that for \emph{finite base fields}, class field theory identifies the maximal abelian quotient with the \(p\)-part of the class group, while nonabelian groups are known in special families only. In the number-field context this must be checked as a statement about finite-degree number-field bases, not silently read as a statement about finite fields of positive characteristic. For finite number-field \(L\), ordinary versus narrow class groups also depends on the infinite-place convention. The triage supplies no precise infinite-base theorem or special-family hypotheses. Its dated general-status claim requires independent verification.
\subsection{Short English statement}
Describe the maximal unramified pro-\(p\) Galois group over the given possibly infinite number-field extension and its continuous abelianization. Ramification conventions, especially at infinity and over an infinite base, matter. The class-field-theory abelian case does not supply the nonabelian structure.
\subsection{Sources}
\begin{itemize}\item Local UnsolvedMath registry, permanent ID 30000488, code catalogue entry; \emph{Unramified Pro-\(p\) Galois Groups over Infinite Extensions}, Oberwolfach Reports Open Problems.\item \emph{Pro-\(p\) Extensions of Global Fields and pro-\(p\) Groups (2007)}, Oberwolfach Reports, DOI 10.4171/owr/2006/25, pp. 1463--1536. \url{https://doi.org/10.4171/owr/2006/25}. Proposed year and year tag: 2006. The source title's escaped dollar delimiters are rendered as mathematics here. No proposer named.\item Supplied literature label \emph{original OWR report}: \url{https://doi.org/10.4171/owr/2006/25}.\item Supplied literature label \emph{Sharifi on unramified pro-\(p\) extensions}: \url{https://arxiv.org/abs/0801.1360}.\item Literature triage checked 2026-08-21, status \texttt{partially\_solved}; record updated 2026-09-16. The \texttt{corrected\_verified} field claims editorial and notation cleanup only. Verify dated claims independently.\end{itemize}
\end{document}
